% !TeX spellcheck = nb_NO
% Kommentaren ovenfor gir instruks til editoren din om at den skal bruke norsk stavekontroll. Dette fungerer bl.a. i TeXstudio.

\documentclass[a4paper,norsk]{article}
% Norsk orddeling
\usepackage[norsk]{babel}
% AMS-pakkene er nyttige
\usepackage{amsmath,amsfonts,amssymb,amsthm}
% For å definere lister, som f.eks. deloppgaver
\usepackage{enumitem}
% Inneholder mye nyttig, slik som \coloneqq
\usepackage{mathtools}
% Penere kalligrafiske fonter
\usepackage[scaled=1.1]{BOONDOX-calo}

% Bruk bedre symboler for ≥, ≤, epsilon og phi
\renewcommand{\geq}{\geqslant}
\renewcommand{\leq}{\leqslant}
\newcommand{\eps} {\varepsilon}
\renewcommand{\epsilon}{\varepsilon}
\renewcommand{\phi}{\varphi}

% Definer de vanligste mengdene og funksjonene
\newcommand{\R}{\mathbb{R}}
\newcommand{\K}{\mathbb{K}}
\newcommand{\C}{\mathbb{C}}
\newcommand{\N}{\mathbb{N}}
\newcommand{\Z}{\mathbb{Z}}
\newcommand{\Q}{\mathbb{Q}}
% \L er definert som noe annet fra før av, så vi må redefinere \L
\renewcommand{\L}{\mathcal{L}}
% Rommet av polynomer
\newcommand{\Pol}{{\mathcal{P}}}
\DeclareMathOperator{\id}{id}
\DeclareMathOperator{\Image}{im}
\DeclareMathOperator{\Kernel}{ker}
\DeclareMathOperator{\Span}{span}
\DeclareMathOperator{\Rank}{rank}
\DeclareMathOperator{\Diag}{diag}
\DeclareMathOperator{\Proj}{proj}
\newcommand{\basis}{{\mathcal{B}}}
\newcommand{\basiss}{{\mathcal{C}}}
\newcommand{\basisss}{{\mathcal{D}}}

% $f \from A \to B$ definerer en funksjon f fra A til B
\newcommand{\from}{\colon}

% \ip{a}{b} er indreproduktet mellom a og b
\newcommand{\ip}[2]{\langle #1, #2\rangle}

% Transponert
\newcommand{\transp}{{\mathsf{T}}}

% Vertikale klammer for absoluttverdier kan gi feil mellomrom, så bruk \abs for absoluttverdier.
\newcommand{\abs}[1]{\mathopen|#1\mathclose|}

% Definer oppgavemiljøet
\newenvironment{oppgave}[1]%
	{\trivlist{}\item\relax\textbf{Løsning på #1.}\medspace}%
	{\endtrivlist}
% Deloppgavelister
\newlist{deloppgave}{enumerate}{1}
\setlist[deloppgave,1]{label=(\textbf{\alph*}),align=left}




\title{Oblig 2 \\ MAT1125, høst 2025}
\author{Ulrik Skre Fjordholm}
\begin{document}
\maketitle


\begin{oppgave}{2.5.2}
La $U$ og $V$ være $n$-dimensjonale vektorrom over $\K$. Av
% En tilde mellom to ord setter inn et mellomrom, men forhindrer at et linjeskift settes inn.
Teorem~2.5.1
finnes isomorfier $S\from U\to \K^n$ og $T\from V\to\K^n$. Da er $\phi\coloneqq T^{-1}\circ S$ en isomorfi fra $U$ til $V$, slik at $U\cong V$.

% Én eller flere blanke linjer lager et nytt avsnitt

(Avbildningen $\phi$ er veldefinert fordi $T^{-1}$ eksisterer og avbilder $\K^n$ til $V$; den er lineær fordi komposisjonen av lineære avbildninger er lineær; og den er bijektiv fordi $S$ og $T^{-1}$ er bijektive, og komposisjonen mellom bijektive funksjoner er bijektiv.)
\end{oppgave}


\begin{oppgave}{2.5.3}
Vi kan skrive
\[
% Tenk på matematiske uttrykk som en del av setningen de forekommer i. Derfor skal tegnsetting (som kommaet her) være med i det matematiske uttrykket.
p(t)=2-t+3t^2 = 2p_0(t)+(-1)p_1(t)+3p_2(t),
\]
så $[p]_\basis = \begin{pmatrix*}[r]
2 \\ -1 \\ 3
\end{pmatrix*}$. Vi har $q_0(t)=1$, $q_1(t)=(t-0)=t$ og $q_2(t)=t(t-1)=t^2-t$. Vi forsøker å ``snike inn'' $q_2$ i uttrykket for $p$, deretter $q_1$, og til sist $q_0$:
% Så langt det er mulig, skal matematiske uttrykk som går over flere linjer justeres horisontalt slik at likhetstegn står rett over hverandre. Bruk "&=" og ikke "=&" (sistnevnte gir feil mellomrom etter likhetstegnet).
\begin{align*}
p(t) &= 2-t+3t^2 = 2-t+3\bigl(\underbrace{(t^2-t)}_{=\,q_2(t)}+t\bigr) \\
&= 2-t+3\bigl(q_2(t)+t\bigr)
= 2+2t+3q_2(t) \\
&= 2q_0(t)+2q_1(t)+3q_2(t).
\end{align*}
Altså er $[p]_\basiss = \begin{pmatrix}
2 \\ 2 \\ 3
\end{pmatrix}$.
\end{oppgave}


\begin{oppgave}{2.5.6}
Av Oppgave~1.4.6 vet vi at om $x\in\K^n$ og $y\in\K^m$, er $[x]_\basis=x$ og $[y]_\basiss=y$. Dermed er
\begin{align*}
Ax &= [Ax]_\basiss = [Tx]_\basiss
= [T]_\basiss^\basis[x]_\basis
= [T]_\basiss^\basis x.
\end{align*}
Siden dette gjelder for alle $x\in\K^n$, må nødvendigvis $[T]_\basiss^\basis = A$.
\end{oppgave}


\begin{oppgave}{2.5.9}
La $n\coloneqq\dim U=\dim V$. Funksjonen $T$ er en isomorfi hvis og bare hvis det finnes en $T^{-1}\in\L(V,U)$ slik at $T^{-1}\circ T=\id_U$ og $T\circ T^{-1}=\id_V$. Av Proposisjon~2.5.8 er da
\[
I_n = [\id_U]^\basis_\basis = [T^{-1}\circ T]^\basis_\basis = [T^{-1}]^\basiss_\basis [T]^\basis_\basiss
\]
og
\[
I_n = [\id_V]^\basiss_\basiss = [T\circ T^{-1}]^\basiss_\basiss
= [T]^\basis_\basiss[T^{-1}]^\basiss_\basis
\]
Altså er $[T^{-1}]^\basiss_\basis$ en inversmatrise av $[T]^\basis_\basiss$, så $[T]^\basis_\basiss$ er inverterbar.

Motsatt, om $[T]^\basis_\basiss$ er inverterbar har den en inversmatrise $A\in M_n(\K)$. Definer $S\from V\to U$ som $Sv \coloneqq A[v]_\basiss$. På samme måte som over ser man da at
\[
I_n = [S\circ T]^\basis_\basis = [T\circ S]^\basiss_\basiss.
\]
Altså er $S\circ T = \id_U$ og $T\circ S=\id_V$, så $S = T^{-1}$.
\end{oppgave}


\begin{oppgave}{2.5.10}
\begin{deloppgave}
\item
For en vilkårlig $k$ er
\[
\frac{dp_k}{dt}(t) = \frac{d}{dt}(t^k) = kt^{k-1} = kp_{k-1}(t).
\]
Dermed er
\[
Tp_0 = 0,\quad Tp_1 = p_0, \quad Tp_2 = 2p_1, \quad Tp_3 = 3p_2.
\]
Hver av disse vektorene i $\Pol_2$ har basisrepresentasjon
\[
[Tp_0]_\basiss = 0, \quad
[Tp_1]_\basiss = \begin{pmatrix}1 \\ 0 \\ 0\end{pmatrix}, \quad
[Tp_2]_\basiss = \begin{pmatrix}0 \\ 2 \\ 0\end{pmatrix}, \quad
[Tp_3]_\basiss = \begin{pmatrix}0 \\ 0 \\ 3\end{pmatrix}.
\]
Vi får dermed basisrepresentasjonen
\[
[T]^\basis_\basiss = \begin{pmatrix}
0 & 1 & 0 & 0 \\
0 & 0 & 2 & 0 \\
0 & 0 & 0 & 3
\end{pmatrix}.
\]
\item
Om $\basis = (p_0,\dots,p_n)$ og $\basiss = (p_0,\dots,p_{n-1})$, får vi som før $Tp_k = kp_{k-1}$ for $k=0,\dots,n$. Dermed er
\[
[Tp_0]_\basiss = 0, \qquad
[Tp_k]_\basiss = ke_{k} \qquad \forall\ k=1,\dots,n
\]
(der $e_1,\dots,e_n$ er standardbasisen i $\R^n$; se Seksjon~0.2.8). Vi får da
\[
[T]^\basis_\basiss = \begin{pmatrix}
0 & 1 & 0 & \cdots & 0 \\
0 & 0 & 2 & \cdots & 0 \\
& \vdots & & \ddots & \vdots \\
0 & 0 & 0 & \cdots & n
\end{pmatrix} \in M_{n\times (n+1)}(\R).
\]
\end{deloppgave}
\end{oppgave}


\begin{oppgave}{2.6.4}
Av Proposisjon~2.6.3 er
\[
[\id]_\basiss^\basis = \begin{pmatrix}u_1 & u_2\end{pmatrix}
= \begin{pmatrix*}[r]
1 & -2 \\ 2 & 2
\end{pmatrix*}.
\]
Dermed er
\[
[\id]_\basis^\basiss = \bigl([\id]_\basiss^\basis\bigr)^{-1}
= \frac16\begin{pmatrix*}[r]
2 & 2 \\ -2 & 1
\end{pmatrix*}.
\]
\end{oppgave}


\begin{oppgave}{2.6.5}
Merk først at $q_0(t)=1$, $q_1(t)=t$ og $q_2(t)=t(t-1)$.

Vi finner først $[\id]_\basis^\basiss$. Om $p=\alpha_0q_0+\alpha_1q_1+\alpha_2q_2$ er et vilkårlig polynom uttrykt i Newton-basisen, er
\begin{align*}
p(t) &= \alpha_0 + \alpha_1t+\alpha_2 t(t-1)
= \alpha_0 + (\alpha_1-\alpha_2)t + \alpha_2 t^2 \\
&= \alpha_0p_0(t) + (\alpha_1-\alpha_2)p_1(t) + \alpha_2 p_2(t).
\end{align*}
Dersom $[p]_\basiss = \begin{pmatrix}
\alpha_0 \\ \alpha_1 \\ \alpha_2
\end{pmatrix}$, er altså $[p]_\basis = \begin{pmatrix}
\alpha_0 \\ \alpha_1-\alpha_2 \\ \alpha_2
\end{pmatrix}$. For at $[p]_\basis = [\id]_\basis^\basiss[p]_\basiss$, må da nødvendigvis
\[
[\id]_\basis^\basiss = \begin{pmatrix*}[r]
1 & 0 & 0 \\ 0 & 1 & -1 \\ 0 & 0 & 1
\end{pmatrix*}.
\]
Dermed er
\[
[\id]_\basiss^\basis = \bigl([\id]_\basis^\basiss\bigr)^{-1}
= \begin{pmatrix*}[r]
1 & 0 & 0 \\ 0 & 1 & 1 \\ 0 & 0 & 1
\end{pmatrix*}.
\]
\end{oppgave}


\end{document}